% ===========================================================
% Chapter 7: Value Function Approximation and Deep Reinforcement Learning
% ===========================================================
\chapter[Value Function Approximation]{Value Function Approximation and Deep Reinforcement Learning}
\label{ch:value-approx}

\begin{quote}
\itshape
The power of neural networks is not that they memorise; it is that they generalise.
In reinforcement learning, this distinction separates methods that scale from methods that do not.
\end{quote}

\bigskip

The preceding six chapters have built a rigorous foundation for reinforcement learning using tabular
representations.  In Chapter~\ref{ch:mdp} we formalised the Markov decision process framework and
derived the Bellman equations; in Chapter~\ref{ch:dp} we showed how dynamic programming solves them
exactly given a model; and in Chapters~\ref{ch:mc} and \ref{ch:td} we developed Monte Carlo and
temporal-difference methods that learn from experience without a model.  These algorithms possess
well-understood convergence guarantees and clean mathematical structure.  Yet they all share one
critical assumption: that we can store and update a separate numerical entry for every
state\,---\,or every state--action pair\,---\,in the problem.

This assumption breaks almost immediately in any realistic application.  An Atari game presents
roughly $10^{120\,000}$ distinct screen configurations; a robotic arm moving in continuous joint
space has an uncountably infinite state space; a language model must assign values to every possible
sequence of tokens.  No computer, present or future, can store a table of that size.  Even if
storage were unlimited, the generalisation properties of tabular methods are poor: learning that
one game state is dangerous conveys no information whatsoever about a visually similar state.

The solution is \emph{function approximation}\index{function approximation}: instead of storing one
number per state, we learn a parameterised function $V_\mathbf{w}(s) \approx V^\pi(s)$ whose
parameter count $d$ is vastly smaller than $|\cS|$.  The same set of parameters simultaneously
encodes information about all states, and updating the parameters after visiting one state
automatically changes the predicted values of nearby states.  This is the mechanism of
generalisation.

This chapter develops function approximation systematically, from first principles to the deep
neural network architectures that power modern reinforcement learning.  Section~\ref{sec:tabular-fail}
quantifies the scalability problem.  Section~\ref{sec:linear-vfa} introduces linear function
approximation, the simplest setting in which convergence can be proved rigorously.
Section~\ref{sec:projected-bellman} analyses the projected Bellman equation, which characterises
the fixed point of TD methods in function-approximation space.
Section~\ref{sec:semi-gradient} examines semi-gradient methods and explains precisely why they are
not true gradient methods.  Section~\ref{sec:deadly-triad} confronts the \emph{deadly triad}: the
conjunction of function approximation, bootstrapping, and off-policy data that can cause divergence
even with linear approximators.  Section~\ref{sec:dqn} presents Deep Q-Networks (DQN), the first
algorithm to demonstrate human-level performance across many diverse tasks by stabilising deep RL
training with two key engineering innovations.  Section~\ref{sec:dqn-extensions} surveys important
algorithmic improvements built upon the DQN foundation.  Section~\ref{sec:convergence-vfa} surveys
what convergence theory can and cannot say about nonlinear function approximation.  Finally,
Section~\ref{sec:applications-vfa} discusses practical considerations that bridge theory and
engineering.

% ===========================================================
\section{Why Tabular Methods Fail}
\label{sec:tabular-fail}
% ===========================================================

\subsection{The Curse of Dimensionality}

To understand the necessity of function approximation, we must first quantify how dramatically the
state space grows with problem complexity.  In a tabular MDP with $|\cS|$ states and $|\cA|$
actions, representing the action-value function $Q : \cS \times \cA \to \R$ requires storing
$|\cS| \cdot |\cA|$ floating-point numbers.  For a small gridworld with $100$ cells and $4$ actions,
that is $400$ numbers\,---\,trivial.  But the size of the state space is not a linear function of
the problem's input dimensionality; it is exponential.

\begin{definition}[Curse of dimensionality]
\label{def:curse-dim}
\index{curse of dimensionality}
Let the state be described by $n$ features, each taking $k$ discrete values.  The total number of
states is $k^n$.  The memory and sample complexity of tabular methods therefore grows
\emph{exponentially} in the number of features $n$.  This phenomenon is called the
\textbf{curse of dimensionality}.
\end{definition}

Three concrete domains make this vivid.

\begin{example}[Atari video games]
\label{ex:atari-states}
\index{Atari}
The Atari 2600 produces $210 \times 160$ pixel frames with $128$ colours per pixel.  Mnih
et al.~(2015) preprocess frames to $84 \times 84$ greyscale and stack four consecutive frames to
capture motion, giving an observation tensor with $84 \times 84 \times 4 = 28\,224$ values, each
an integer in $\{0, \ldots, 255\}$.  The number of distinct observations is therefore at most
$256^{28\,224}$, a number with roughly $68\,000$ decimal digits.  A tabular $Q$-function over
this space is not merely impractical; it is conceptually absurd.  Yet DQN, using a convolutional
neural network with $\sim 1.5$ million parameters, achieved superhuman performance on many of these
games.  The network succeeds precisely because it generalises across visually similar frames.
\end{example}

\begin{example}[Continuous robotics]
\label{ex:robotics-states}
\index{robotics!state space}
A simulated humanoid robot (e.g., the MuJoCo Humanoid-v4 environment) has a state vector
$s \in \R^{376}$ comprising joint positions, joint velocities, and contact forces.  This state
space is not merely large\,---\,it is continuous and uncountably infinite.  No discretisation
scheme can reduce it to a manageable table without introducing severe approximation errors.
Moreover, the action space is $\R^{17}$, so the $\argmax$ operation central to $Q$-learning
becomes a continuous optimisation problem at every step.
\end{example}

\begin{example}[Language and text]
\label{ex:language-states}
\index{language models!state space}
In applications of RL to language model fine-tuning (discussed further in
Chapter~\ref{ch:practical-rlhf}), the state is the current token prefix\,---\,the entire
conversation history up to the current step.  With a vocabulary of size $50\,000$ and a context
window of $2\,048$ tokens, the number of possible prefixes is astronomically large.  Furthermore,
the semantics of two prefixes differing by a single word can be nearly identical, so any useful
representation must achieve semantic generalisation, not just syntactic.
\end{example}

\subsection{The Generalisation Argument}

Beyond the pure storage problem, tabular methods have a fundamental statistical inefficiency: every
state is treated as a completely separate object.  Learning that the value of state $s$ is high
conveys no information about state $s'$, even if $s$ and $s'$ are visually identical up to the
position of one pixel.  In the language of statistics, tabular methods have zero
\emph{generalisation}: each entry must be estimated independently, from transitions that visit
that specific state.

In large state spaces, most states are visited at most once, or never at all, during any reasonable
training run.  Tabular methods assign such states their initial values (often zero) indefinitely.
Function approximation addresses this by sharing statistical strength across states: a visit to
state $s$ implicitly updates the estimated values of all states $s'$ for which $\varphi(s')$ is
similar to $\varphi(s)$, where $\varphi$ denotes the feature representation.  This
\emph{implicit regularisation} is the core benefit of parameterised value functions.

% ===========================================================
\section{Linear Value Function Approximation}
\label{sec:linear-vfa}
% ===========================================================

\subsection{Feature Vectors and the Linear Ansatz}

The simplest parameterised approximation represents the value function as a linear combination of
fixed features.

\begin{definition}[Linear value function approximation]
\label{def:linear-vfa}
\index{linear value function approximation}
Let $\varphi : \cS \to \R^d$ be a \textbf{feature map}\index{feature map} assigning each state a
feature vector.  A \textbf{linear value function approximation} takes the form
\begin{equation}
\label{eq:linear-vfa}
V_\mathbf{w}(s) = \mathbf{w}^\top \varphi(s) = \sum_{i=1}^d w_i \varphi_i(s),
\end{equation}
where $\mathbf{w} \in \R^d$ is the weight vector.  For action-value functions one uses features
$\varphi(s, a) \in \R^d$ and writes $Q_\mathbf{w}(s, a) = \mathbf{w}^\top \varphi(s, a)$.
\end{definition}

The linear ansatz is attractive for three reasons.  First, it is the simplest possible
parameterisation: the approximation is uniquely determined by $d$ real numbers regardless of how
large $|\cS|$ is.  Second, many interesting feature maps can be designed by hand\,---\,tile
codings, radial basis functions, polynomial bases, Fourier features\,---\,each encoding specific
prior beliefs about the structure of the value function.  Third, and most importantly, linear
function approximation is the setting where TD convergence can be established with mathematical
rigour, as we will see in Section~\ref{sec:projected-bellman}.

\begin{example}[Polynomial features for a one-dimensional corridor]
\label{ex:poly-features}
\index{polynomial features}
Consider a corridor of $N = 1\,000$ discrete positions.  In the tabular case, $V^\pi$ requires
$N = 1\,000$ parameters.  If, however, the true value function is smooth\,---\,which is likely if
the reward is concentrated at one end and the discount factor is moderate\,---\,we can approximate
it with the polynomial basis
\[
\varphi(s) = \bigl(1,\; s/N,\; (s/N)^2,\; (s/N)^3\bigr)^\top \in \R^4.
\]
This reduces the parameter count from $1\,000$ to $4$, at the cost of being unable to represent
rapidly varying value functions exactly.  The approximation error incurred by restricting to this
function class is called the \textbf{representation error}\index{representation error}, and it is
unavoidable whenever the true value function lies outside the approximation family.
\end{example}

\subsection{Tile Coding}

A particularly useful feature map for continuous state spaces is \emph{tile coding}%
\index{tile coding}, developed by Sutton~(1996).  The state space $\R^n$ is covered by multiple
overlapping grids, called \emph{tilings}.  For a given state $s$, exactly one tile in each tiling
is active, giving a binary feature vector whose non-zero entries indicate which tile is active in
each tiling.  If there are $m$ tilings each partitioning the space into $k$ tiles, the feature
vector has dimension $mk$ and contains exactly $m$ ones.  The resulting linear approximation is
equivalent to a lookup table of tile values, but the overlapping structure of the tilings provides
smooth interpolation between nearby states.

\subsection{Linear TD(0)}
\label{subsec:linear-td}

Recall from Chapter~\ref{ch:td} that TD(0) for policy evaluation maintains a value estimate and
updates it after each transition $(S_t, A_t, R_{t+1}, S_{t+1})$ using the TD error
\[
\delta_t = R_{t+1} + \gamma V(S_{t+1}) - V(S_t).
\]
When $V$ is parameterised as $V_\mathbf{w}(s) = \mathbf{w}^\top \varphi(s)$, the natural
extension is to update the weight vector in the direction that reduces the squared TD error.  This
yields the \textbf{semi-gradient TD(0)}\index{semi-gradient TD(0)} update (analysed in detail in
Section~\ref{sec:semi-gradient}):
\begin{equation}
\label{eq:linear-td-update}
\mathbf{w} \leftarrow \mathbf{w} + \alpha\, \delta_t\, \varphi(S_t),
\end{equation}
where $\delta_t = R_{t+1} + \gamma \mathbf{w}^\top \varphi(S_{t+1}) - \mathbf{w}^\top \varphi(S_t)$.
Observe that in the linear case $\nabla_\mathbf{w} V_\mathbf{w}(S_t) = \varphi(S_t)$, so
\eqref{eq:linear-td-update} is the semi-gradient step applied to $V_\mathbf{w}$.

A critical property of this update is that changing $\mathbf{w}$ affects \emph{all states
simultaneously}: the value prediction at state $s'$ changes by $\alpha \delta_t \, \varphi(s')
\cdot \varphi(S_t)$, proportional to how similar $s'$ is to $S_t$ under the feature inner product.
This is the mechanism of generalisation in linear function approximation.

\begin{theorem}[Convergence of linear TD(0) with on-policy sampling]
\label{thm:linear-td-convergence}
\index{linear TD(0)!convergence}
Let $\pi$ be a fixed policy and let $\mu$ be its stationary distribution over $\cS$ (assumed
unique and positive everywhere).  Assume the feature matrix $\Phi \in \R^{|\cS| \times d}$ has
full column rank.  If samples $(S_t, R_{t+1}, S_{t+1})$ are drawn on-policy, i.e., $S_t \sim \mu$
in steady state, and the step sizes satisfy $\sum_t \alpha_t = \infty$ and
$\sum_t \alpha_t^2 < \infty$, then the linear TD(0) iterates $\mathbf{w}_t$ converge
almost surely to a unique weight vector $\mathbf{w}_{\mathrm{TD}}$ satisfying the projected
Bellman equation $\Phi \mathbf{w}_{\mathrm{TD}} = \Pi_\mu T^\pi (\Phi \mathbf{w}_{\mathrm{TD}})$
(defined in Section~\ref{sec:projected-bellman}).
\end{theorem}

The proof of Theorem~\ref{thm:linear-td-convergence} relies on the theory of stochastic
approximation and the contraction properties of the projected Bellman operator, which we now
develop.

% ===========================================================
\section{The Projected Bellman Equation}
\label{sec:projected-bellman}
% ===========================================================

\subsection{Geometry of Function Approximation}

To understand what linear TD(0) actually converges to, we need to think geometrically.  The linear
function class $\mathcal{F} = \{\Phi \mathbf{w} : \mathbf{w} \in \R^d\}$ is a $d$-dimensional
subspace of $\R^{|\cS|}$ (identifying value functions with vectors indexed by states).  The
Bellman operator $T^\pi$ maps any value function $v \in \R^{|\cS|}$ to a new function
$T^\pi v$, but in general $T^\pi v \notin \mathcal{F}$ even when $v \in \mathcal{F}$.
Consequently, TD methods operating within $\mathcal{F}$ cannot follow $T^\pi$ directly.

The resolution is to \emph{project} $T^\pi v$ back onto $\mathcal{F}$ after each Bellman step.

\begin{definition}[$\mu$-weighted inner product and projection]
\label{def:mu-projection}
\index{projection operator!$\mu$-weighted}
Let $\mu : \cS \to (0, \infty)$ be a positive weighting (typically the on-policy stationary
distribution of $\pi$).  Define the \textbf{$\mu$-weighted norm}\index{$\mu$-norm}
\begin{equation}
\label{eq:mu-norm}
\|v\|_\mu^2 = \sum_{s \in \cS} \mu(s)\, v(s)^2 = \mathbf{v}^\top D_\mu \mathbf{v},
\end{equation}
where $D_\mu = \mathrm{diag}(\mu)$.  The \textbf{orthogonal projection} of $v$ onto
$\mathcal{F}$ in this norm is
\begin{equation}
\label{eq:projection-def}
\Pi_\mu v = \Phi\bigl(\Phi^\top D_\mu \Phi\bigr)^{-1} \Phi^\top D_\mu v
= \argmin_{f \in \mathcal{F}} \|f - v\|_\mu.
\end{equation}
\end{definition}

The matrix $\Pi_\mu = \Phi(\Phi^\top D_\mu \Phi)^{-1}\Phi^\top D_\mu$ is the orthogonal
projector in the $\mu$-weighted inner product.  It is idempotent ($\Pi_\mu^2 = \Pi_\mu$) and
satisfies $\|I - \Pi_\mu\|_\mu \leq 1$, meaning that projection can only bring a vector
\emph{closer} to $\mathcal{F}$, never push it further away.

\begin{definition}[Projected Bellman equation]
\label{def:projected-bellman}
\index{projected Bellman equation}
The \textbf{projected Bellman equation} for policy $\pi$ is
\begin{equation}
\label{eq:projected-bellman}
\hat{V} = \Pi_\mu T^\pi \hat{V}, \qquad \hat{V} \in \mathcal{F}.
\end{equation}
A solution $\hat{V} = \Phi \mathbf{w}_{\mathrm{TD}}$ is called the
\textbf{TD fixed point}\index{TD fixed point}.
\end{definition}

The projected Bellman equation has a clear geometric interpretation: find a function in the
approximation class $\mathcal{F}$ that, after one Bellman backup and re-projection, returns to
itself.  It is the best we can hope for when exact Bellman consistency is unachievable within
$\mathcal{F}$.

\begin{example}[Projected Bellman equation with three states]
\label{ex:three-state-projection}
Suppose $|\cS| = 3$ with uniform weighting $\mu = (1/3, 1/3, 1/3)$, and our feature matrix is
\[
\Phi = \begin{pmatrix} 1 & 0 \\ 0 & 1 \\ 1 & 1 \end{pmatrix} \in \R^{3 \times 2}.
\]
The approximation family $\mathcal{F}$ has dimension $2$, so values at all three states cannot be
set independently: $V(s_3) = V(s_1) + V(s_2)$ always.  The projection $\Pi_\mu v$ of any
$v \in \R^3$ finds the weights $\mathbf{w}^* = \argmin_\mathbf{w} \|v - \Phi \mathbf{w}\|_\mu^2$
by solving $\Phi^\top D_\mu \Phi \mathbf{w}^* = \Phi^\top D_\mu v$, yielding
\[
\mathbf{w}^* = \tfrac{1}{2}\begin{pmatrix} v_1 + v_3 - v_2 \\ v_2 + v_3 - v_1 \end{pmatrix},
\]
so the projection ``splits the contribution of $s_3$'' across the two basis directions.
\end{example}

\subsection{Existence of the TD Fixed Point}

\begin{theorem}[TD fixed point existence and uniqueness]
\label{thm:td-fixed-point}
\index{TD fixed point!existence}
Under the assumptions of Theorem~\ref{thm:linear-td-convergence}, the projected Bellman operator
$\Pi_\mu T^\pi : \mathcal{F} \to \mathcal{F}$ is a contraction in the $\mu$-norm with
contraction factor $\gamma$:
\begin{equation}
\label{eq:projected-contraction}
\|\Pi_\mu T^\pi u - \Pi_\mu T^\pi v\|_\mu \leq \gamma \|u - v\|_\mu
\qquad \forall\, u, v \in \mathcal{F}.
\end{equation}
Consequently, the projected Bellman equation~\eqref{eq:projected-bellman} has a unique solution
$\hat{V} = \Phi \mathbf{w}_{\mathrm{TD}} \in \mathcal{F}$.
\end{theorem}

\begin{proof}
The Bellman operator $T^\pi$ is itself a contraction in the $\mu$-norm with factor $\gamma$
(this follows from the contraction property of $T^\pi$ in the sup-norm and the equivalence of
these norms for positive $\mu$).  Formally, for any $u, v$:
\[
\|T^\pi u - T^\pi v\|_\mu \leq \gamma \|u - v\|_\mu.
\]
Since $\Pi_\mu$ is a non-expansive projection ($\|\Pi_\mu x\|_\mu \leq \|x\|_\mu$ for all $x$),
we have
\[
\|\Pi_\mu T^\pi u - \Pi_\mu T^\pi v\|_\mu
\leq \|T^\pi u - T^\pi v\|_\mu
\leq \gamma \|u - v\|_\mu.
\]
Because $\gamma < 1$ and $\mathcal{F}$ is a complete metric space (a closed subspace of
$\R^{|\cS|}$), the Banach fixed-point theorem guarantees a unique fixed point.
\end{proof}

\subsection{Approximation Error at the TD Fixed Point}

The TD fixed point $\hat{V}$ minimises a particular measure of error, but it does not minimise the
mean-squared error $\|\hat{V} - V^\pi\|_\mu^2$ directly.  The following classical bound
(Tsitsiklis and Van Roy, 1997) quantifies how far the TD fixed point deviates from the best
possible approximation.

\begin{theorem}[TD fixed-point approximation bound]
\label{thm:td-bound}
\index{TD fixed point!approximation error}
Let $\hat{V}^* = \Pi_\mu V^\pi$ be the best approximation of $V^\pi$ in $\mathcal{F}$.  Then
\begin{equation}
\label{eq:td-bound}
\|\hat{V} - V^\pi\|_\mu \leq \frac{1}{\sqrt{1 - \gamma^2}}\, \|\hat{V}^* - V^\pi\|_\mu.
\end{equation}
\end{theorem}

The factor $1/\sqrt{1-\gamma^2}$ is a ``price'' paid for bootstrapping: the TD fixed point may be
up to $1/\sqrt{1-\gamma^2}$ times worse than the best-in-class approximation.  For $\gamma = 0.99$
this factor is approximately $7.1$, illustrating that bootstrapping can significantly amplify
approximation errors in problems with long time horizons.

% ===========================================================
\section{Semi-Gradient Methods}
\label{sec:semi-gradient}
% ===========================================================

\subsection{True Gradient Descent for Prediction}

Before examining semi-gradient methods, it is instructive to understand what true gradient descent
would look like for the prediction problem.  Define the \textbf{mean-squared value
error}\index{mean-squared value error}
\begin{equation}
\label{eq:msve}
\overline{\mathrm{VE}}(\mathbf{w})
= \|V^\pi - V_{\mathbf{w}}\|_{\mu}^{2}
= \sum_{s \in \cS} \mu(s)\bigl[V^\pi(s) - V_\mathbf{w}(s)\bigr]^2,
\end{equation}
where $\mu$ is the on-policy distribution.  The true gradient of $\overline{\mathrm{VE}}$
with respect to $\mathbf{w}$ is
\begin{equation}
\label{eq:true-grad-msve}
-\frac{1}{2}\nabla_\mathbf{w}\overline{\mathrm{VE}}(\mathbf{w})
= \sum_s \mu(s)\bigl[V^\pi(s) - V_\mathbf{w}(s)\bigr]\nabla_\mathbf{w} V_\mathbf{w}(s).
\end{equation}
Unfortunately, $V^\pi(s)$ is unknown, so we cannot evaluate this gradient directly.  Monte Carlo
methods substitute the sample return $G_t \approx V^\pi(S_t)$, giving the update
\[
\mathbf{w} \leftarrow \mathbf{w} + \alpha\bigl[G_t - V_\mathbf{w}(S_t)\bigr]\nabla_\mathbf{w} V_\mathbf{w}(S_t).
\]
This is a true stochastic gradient step because $\E[G_t \mid S_t] = V^\pi(S_t)$ and the
gradient $\nabla_\mathbf{w} V_\mathbf{w}(S_t)$ does not appear in the target $G_t$.

\subsection{Why TD Breaks True Gradient Descent}

TD methods replace $G_t$ with the one-step bootstrap target $R_{t+1} + \gamma V_\mathbf{w}(S_{t+1})$.
The resulting update is
\begin{equation}
\label{eq:semi-grad-td}
\mathbf{w} \leftarrow \mathbf{w} + \alpha\underbrace{\bigl[R_{t+1} + \gamma V_\mathbf{w}(S_{t+1}) - V_\mathbf{w}(S_t)\bigr]}_{\delta_t}\nabla_\mathbf{w} V_\mathbf{w}(S_t).
\end{equation}
This \emph{looks} like a gradient step, but it is not.

\begin{definition}[Semi-gradient method]
\label{def:semi-gradient}
\index{semi-gradient method}
A \textbf{semi-gradient} method is one in which the update direction is computed as if $\delta_t$
were an error signal independent of $\mathbf{w}$, i.e., the gradient is taken only through
$V_\mathbf{w}(S_t)$ and not through the target $V_\mathbf{w}(S_{t+1})$.  The update
\eqref{eq:semi-grad-td} is a semi-gradient step because the gradient of $\delta_t$ with respect
to $\mathbf{w}$ via its dependence on $V_\mathbf{w}(S_{t+1})$ is ignored.
\end{definition}

To see why the gradient through $V_\mathbf{w}(S_{t+1})$ is omitted, note that if we computed
the true gradient of $\frac{1}{2}\delta_t^2$ we would obtain
\[
\nabla_\mathbf{w}\tfrac{1}{2}\delta_t^2
= \delta_t\bigl(\gamma\nabla_\mathbf{w}V_\mathbf{w}(S_{t+1}) - \nabla_\mathbf{w}V_\mathbf{w}(S_t)\bigr).
\]
The TD semi-gradient drops the $\gamma\nabla_\mathbf{w}V_\mathbf{w}(S_{t+1})$ term.  This
omission means the update direction is \emph{not} the negative gradient of any fixed objective
function.  As a consequence, standard convergence theorems for stochastic gradient descent do not
apply directly.

\begin{remark}
\label{rem:semi-grad-motivation}
The omission of the bootstrap gradient is deliberate, not an error.  If one were to include it,
the resulting \emph{residual gradient} algorithm converges in the mean-squared Bellman error but
suffers from a ``double-sampling'' problem: computing the full gradient requires two independent
samples of $S_{t+1}$ for the same $(S_t, A_t)$, which is typically unavailable in a single
interaction stream.  Semi-gradient methods avoid this by treating the target as fixed; their
practical convergence properties are generally superior despite the lack of a gradient guarantee.
\end{remark}

\subsection{Semi-Gradient TD Algorithm}

\begin{algorithm}[ht]
\caption{Semi-Gradient TD(0) for Policy Evaluation}
\label{alg:semi-grad-td0}
\KwIn{Policy $\pi$, feature map $\varphi$, step size $\alpha > 0$, discount $\gamma \in [0,1)$}
Initialise $\mathbf{w} \in \R^d$ (e.g., $\mathbf{w} \leftarrow \mathbf{0}$)\;
\For{each episode}{
  Initialise state $S_0$\;
  \For{$t = 0, 1, 2, \ldots$ until $S_t$ is terminal}{
    Choose $A_t \sim \pi(\cdot \mid S_t)$\;
    Observe $R_{t+1}$ and $S_{t+1}$\;
    $\delta_t \leftarrow R_{t+1} + \gamma \mathbf{w}^\top \varphi(S_{t+1}) - \mathbf{w}^\top \varphi(S_t)$\;
    $\mathbf{w} \leftarrow \mathbf{w} + \alpha\, \delta_t\, \varphi(S_t)$\;
  }
}
\end{algorithm}

\subsection{Convergence of Semi-Gradient TD}

For the linear case, semi-gradient TD(0) under on-policy sampling converges to the TD fixed point
$\mathbf{w}_{\mathrm{TD}}$ characterised by Theorem~\ref{thm:td-fixed-point}.  The argument
relies on the \emph{ordinary differential equation (ODE) method}\index{ODE method}: one shows that
the mean update field
\[
\bar{f}(\mathbf{w}) = \E_\mu[\delta_t \varphi(S_t) \mid \mathbf{w}]
= \Phi^\top D_\mu (T^\pi(\Phi\mathbf{w}) - \Phi\mathbf{w})
\]
is a contraction toward $\mathbf{w}_{\mathrm{TD}}$, and then appeals to a standard theorem
(Kushner and Yin, 2003) that says stochastic approximation converges to the equilibrium of the
associated ODE.

For \emph{nonlinear} function approximators, even on-policy semi-gradient TD is not guaranteed to
converge.  In practice, neural-network-based TD methods often converge to useful solutions, but
counter-examples exist.  Convergence in the nonlinear case remains an active research area.

% ===========================================================
\section{The Deadly Triad}
\label{sec:deadly-triad}
% ===========================================================

\subsection{Three Ingredients for Instability}

The convergence guarantee of Theorem~\ref{thm:linear-td-convergence} hinges critically on the
words ``on-policy sampling''.  When any one of three conditions is violated, the well-behaved
landscape of linear TD breaks down.  When all three violations occur simultaneously, even linear
function approximation can diverge.  Sutton and Barto~(2018) term this dangerous combination the
\textbf{deadly triad}\index{deadly triad}.

\begin{definition}[The deadly triad]
\label{def:deadly-triad}
\index{deadly triad}
The \textbf{deadly triad} refers to the simultaneous presence of three elements:
\begin{enumerate}
\item \textbf{Function approximation:} using a parameterised function
  $V_\mathbf{w}$ (or $Q_\mathbf{w}$) with $d \ll |\cS|$ instead of a separate table entry per
  state.
\item \textbf{Bootstrapping:} forming targets from the current estimate, as in TD methods, rather
  than waiting for the full return, as in Monte Carlo.
\item \textbf{Off-policy training:} updating parameters using data generated by a \emph{behaviour
  policy} $b$ that differs from the \emph{target policy} $\pi$ being evaluated or improved.
\end{enumerate}
The combination of all three can cause parameter divergence even with linear approximators and
even in the limit of infinite data.
\end{definition}

\begin{figure}[ht]
\centering
\begin{tikzpicture}[
  every node/.style={font=\small},
  tri/.style={draw, fill=gray!15, regular polygon, regular polygon sides=3,
              minimum size=5.5cm, inner sep=0pt},
  label/.style={font=\small\bfseries, align=center}
]
  % Deadly triad triangle
  \node[tri] (T) {};
  % Vertices
  \coordinate (A) at (T.corner 1);  % top
  \coordinate (B) at (T.corner 2);  % bottom-left
  \coordinate (C) at (T.corner 3);  % bottom-right

  \node[label, above=0.45cm of A, align=center] {Function\\Approximation};
  \node[label, below=0.25cm of B, align=right] {Bootstrapping};
  \node[label, below=0.25cm of C, align=left] {Off-Policy\\Training};

  \node[font=\itshape\small, align=center] at ($(A)!0.5!(B)!0.33!(C)$)
    {Potential\\Divergence};

  % Safe-zone annotations on edges (sloped along each edge)
  \node[font=\tiny, sloped, above, inner sep=1pt] at ($(A)!0.5!(B)$) {MC (no bootstrap)};
  \node[font=\tiny, above, inner sep=2pt] at ($(B)!0.5!(C)+(0,0.1)$) {Tabular (no approx.)};
  \node[font=\tiny, sloped, above, inner sep=1pt] at ($(A)!0.5!(C)$) {On-policy TD};
\end{tikzpicture}
\caption{The deadly triad.  Each edge of the triangle is ``safe'': two of the three ingredients
together are manageable.  Only when all three vertices are active simultaneously does instability
become unavoidable in general.  The edge labels indicate which element is absent along each side.}
\label{fig:deadly-triad}
\end{figure}

It is important to understand that each pair of ingredients is manageable:
\begin{itemize}[leftmargin=*]
\item \emph{Function approximation + bootstrapping, on-policy:} Linear TD(0) converges to the TD
  fixed point (Theorem~\ref{thm:linear-td-convergence}).
\item \emph{Function approximation + off-policy, no bootstrap:} Off-policy Monte Carlo with linear
  function approximation converges because the gradient is a true gradient of the MSVE.
\item \emph{Bootstrapping + off-policy, tabular:} Q-learning converges to $Q^*$ in the tabular
  setting (Chapter~\ref{ch:td}, Theorem on Q-learning convergence).
\end{itemize}

The problem emerges only when all three are combined.

\subsection{Baird's Counterexample}
\label{subsec:baird}

The clearest illustration of the deadly triad is \emph{Baird's counterexample}\index{Baird's counterexample}
(Baird, 1995).  Consider a simple MDP with $7$ states and $2$ actions.  Six states are arranged in
a star pattern; the seventh is a central hub.  The behaviour policy always chooses the first action
(which transitions to one of the six star states, uniformly at random) with probability $\frac{6}{7}$
and the second action (which transitions to the hub) with probability $\frac{1}{7}$.
The target policy always takes the second action.  All rewards are zero, and $\gamma = 0.99$.

The feature vectors are chosen so that the approximation family is linear.  Specifically, the
feature vector for hub state $s_7$ is $(0,0,0,0,0,0,1,2)^\top$ and for star state $s_i$
($i = 1, \ldots, 6$) is the vector with a $2$ in position $i$ and a $1$ in position $8$.  The
weight vector $\mathbf{w} \in \R^8$ is initialised to $(1,1,1,1,1,1,10,1)^\top$.

\begin{theorem}[Baird's counterexample]
\label{thm:baird}
\index{Baird's counterexample!divergence}
Under semi-gradient off-policy TD(0) with the feature vectors described above, the weight vector
$\mathbf{w}_t$ diverges to infinity almost surely, even though the feature matrix has full column
rank and the step size satisfies the standard Robbins--Monro conditions.
\end{theorem}

The mechanism of divergence is subtle: the off-policy distribution $\mu$ concentrates too much
weight on the star states (which are frequently visited under the behaviour policy) while the
Bellman updates primarily ``push'' the values of those states up, since the target policy prefers
to go to the hub.  The projection back onto the function class then amplifies this push rather
than damping it, creating a positive feedback loop.

\begin{remark}[What Baird's example teaches us]
\label{rem:baird-lesson}
Baird's example is not a pathological curiosity; it is a fundamental result about the limits of
TD methods.  It tells us that the combination of semi-gradient bootstrapping and off-policy data
is inherently dangerous, regardless of how well-behaved the feature representation is.  The
solutions explored in practice\,---\,importance sampling corrections, gradient-TD methods
(TDC, GTD2), and the engineering stabilisations of DQN\,---\,all represent different strategies
for escaping the triad's dangers without sacrificing the benefits of off-policy data.
\end{remark}

% ===========================================================
\section{Deep Q-Networks}
\label{sec:dqn}
% ===========================================================

\subsection{From Linear to Nonlinear Approximation}

The successes of linear function approximation are real but limited: the method requires
hand-designed features, and in high-dimensional perceptual domains such as raw image pixels, no
human designer can specify the right features in advance.  The natural next step is to use a
deep neural network\index{deep neural network} as the approximator, letting the network learn its
own internal representation from data.

The challenge is that deep networks bring all the difficulties of the deadly triad into sharp
focus.  In Q-learning with a neural network $Q_\theta$ and greedy action selection, the target for
state--action pair $(S_t, A_t)$ is
\begin{equation}
\label{eq:naive-deep-q-target}
y_t = R_{t+1} + \gamma \max_{a'} Q_\theta(S_{t+1}, a'),
\end{equation}
and the loss minimised is $(y_t - Q_\theta(S_t, A_t))^2$.  But both the prediction $Q_\theta(S_t, A_t)$
and the target $y_t$ depend on the same parameters $\theta$, so each gradient step simultaneously
changes the prediction \emph{and} the target.  This creates a moving-target problem: the network
chases a quantity that shifts with every update, leading to oscillations and divergence in practice.

Two further pathologies compound the problem.  The transitions $(S_t, A_t, R_{t+1}, S_{t+1})$
collected by a sequential agent are \emph{highly correlated}: consecutive transitions share the
same state sequence.  Training a neural network on correlated mini-batches violates the i.i.d.\
assumption of stochastic gradient descent and causes the network to overfit to recent experience,
``forgetting'' earlier learning.  Additionally, the distribution of training data changes as the
policy improves, making the optimisation landscape non-stationary.

\subsection{Experience Replay}
\label{subsec:experience-replay}

The first stabilising innovation introduced by Mnih et al.~(2013, 2015) is
\textbf{experience replay}\index{experience replay}.  As the agent interacts with the environment,
each transition $(S_t, A_t, R_{t+1}, S_{t+1})$ is stored in a \emph{replay buffer}\index{replay buffer}
$\mathcal{D}$ of capacity $N$.  During training, mini-batches of transitions are sampled
\emph{uniformly at random} from $\mathcal{D}$:
\begin{equation}
\label{eq:replay-sampling}
\{(s_i, a_i, r_i, s'_i)\}_{i=1}^B \sim \mathrm{Uniform}(\mathcal{D}).
\end{equation}
This simple change has two important effects.  First, it \emph{decorrelates} training data:
successive mini-batches are drawn from many different time steps, breaking the temporal
correlation that destabilises on-line updates.  Second, each transition is used for potentially
many gradient steps, improving data efficiency.  The cost is that replay requires storing a large
buffer in memory and introduces an implicit off-policy component: transitions stored long ago were
generated by an older, different policy.

\subsection{Target Networks}
\label{subsec:target-networks}

The second stabilising innovation is the \textbf{target network}\index{target network}.  DQN
maintains two networks with the same architecture: the \emph{online network} $Q_\theta$ (also
called the \emph{behaviour network}) and the \emph{target network} $Q_{\theta^-}$.  The online
network is updated at every gradient step.  The target network's parameters $\theta^-$ are held
fixed for $C$ steps (typically $C = 10\,000$) and then synchronised with the online network:
$\theta^- \leftarrow \theta$.

With the target network, the loss function becomes
\begin{equation}
\label{eq:dqn-loss}
L(\theta) = \E_{(s,a,r,s') \sim \mathcal{D}}\!\left[
\Bigl(r + \gamma \max_{a'} Q_{\theta^-}(s', a') - Q_\theta(s, a)\Bigr)^2
\right].
\end{equation}
The gradient is taken only with respect to $\theta$, treating $\theta^-$ as a constant.  By
freezing the target for $C$ steps, the method reduces the correlation between the parameters
being updated and the target they are chasing.  The target is not perfectly stationary\,---\,it
still moves every $C$ steps\,---\,but in practice this significantly stabilises training.


\begin{figure}[ht]
\centering
\begin{tikzpicture}[scale=0.9, transform shape,
  box/.style={draw, rounded corners, minimum width=2.8cm, minimum height=0.9cm,
              font=\small, fill=gray!10, align=center},
  buf/.style={draw, rounded corners, minimum width=2.8cm, minimum height=0.9cm,
              font=\small, fill=blue!10, align=center},
  lossbox/.style={draw, rounded corners, minimum width=2.8cm, minimum height=0.9cm,
              font=\small, fill=orange!10, align=center},
  arr/.style={-{Latex[length=2mm]}, thick},
  darr/.style={-{Latex[length=2mm]}, thick, dashed}
]

  % Row 1: Online Network (left) and Target Network (right)
  \node[box] (online) at (0,0) {Online Network\\$Q_\theta$};
  \node[box] (target) at (5.5,0) {Target Network\\$Q_{\theta^-}$};

  % Row 2: Loss centred between the two networks
  \node[lossbox] (loss) at (2.75,-2.8) {Loss $L(\theta)$};

  % Row 3: Replay Buffer centred below Loss
  \node[buf] (buf) at (2.75,-5.0) {Replay Buffer $\mathcal{D}$};

  % Dashed arrow: copy weights
  \draw[darr] (online.north) to[bend left=20]
    node[above, font=\scriptsize]{copy every $C$ steps} (target.north);

  % Online Network -> Loss
  \draw[arr] (online.south) -- (loss.north west)
    node[midway, left, font=\scriptsize, xshift=-2pt]{predict $Q_\theta(s,a)$};

  % Target Network -> Loss
  \draw[arr] (target.south) -- (loss.north east)
    node[midway, right, font=\scriptsize, xshift=2pt]{target $\max_{a'}Q_{\theta^-\!}(s',a')$};

  % Replay Buffer -> Loss
  \draw[arr] (buf.north) -- (loss.south)
    node[midway, right, font=\scriptsize, xshift=2pt]{mini-batch};

  % Loss -> Online Network (gradient, curved back on the far left)
  \draw[arr] (loss.west) to[out=180, in=250, looseness=1.2]
    node[left, font=\scriptsize, xshift=-4pt]{$\nabla_\theta L$}
    (online.south west);

\end{tikzpicture}
\caption{DQN architecture.  The online network $Q_\theta$ selects actions and is updated every step.
The target network $Q_{\theta^-}$ computes the TD target and is synchronised every $C$ steps.
Transitions are stored in the replay buffer $\mathcal{D}$ and sampled in random mini-batches for
training.  Dashed arrow: periodic hard copy $\theta^- \leftarrow \theta$.}
\label{fig:dqn-arch}
\end{figure}

% \begin{figure}[ht]
% \centering
% \begin{tikzpicture}[scale=0.9, transform shape,
%   box/.style={draw, rounded corners, minimum width=2.8cm, minimum height=0.9cm,
%               font=\small, fill=gray!10, align=center},
%   buf/.style={draw, rounded corners, minimum width=2.8cm, minimum height=0.9cm,
%               font=\small, fill=blue!10, align=center},
%   lossbox/.style={draw, rounded corners, minimum width=2.8cm, minimum height=0.9cm,
%               font=\small, fill=orange!10, align=center},
%   arr/.style={-{Latex[length=2mm]}, thick},
%   darr/.style={-{Latex[length=2mm]}, thick, dashed}
% ]
%   % Row 1: Online Network (left) and Target Network (right)
%   \node[box] (online) {Online Network\\$Q_\theta$};
%   \node[box, right=3.8cm of online] (target) {Target Network\\$Q_{\theta^-}$};

%   % Row 2: Loss centred below
%   \node[lossbox, below=2.2cm of $(online.south)!0.5!(target.south)$] (loss) {Loss $L(\theta)$};

%   % Row 3: Replay Buffer centred below Loss
%   \node[buf, below=1.4cm of loss] (buf) {Replay Buffer $\mathcal{D}$};

%   % Dashed arrow: copy weights
%   \draw[darr] (online.north) to[bend left=18]
%     node[above, font=\scriptsize]{copy every $C$ steps} (target.north);

%   % Online Network -> Loss (straight down from online to the loss node)
%   \draw[arr] (online.south) -- (online.south |- loss.north)
%     node[midway, right, font=\scriptsize, xshift=2pt]{predict $Q_\theta(s,a)$};

%   % Target Network -> Loss (straight down from target to the loss node)
%   \draw[arr] (target.south) -- (target.south |- loss.north)
%     node[midway, left, font=\scriptsize, xshift=-2pt]{target $\max_{a'}Q_{\theta^-\!}(s',a')$};

%   % Replay Buffer -> Loss (straight up)
%   \draw[arr] (buf.north) -- node[right, font=\scriptsize, xshift=2pt]{mini-batch} (loss.south);

%   % Loss -> Online Network (curved back on the far left)
%   \draw[arr] (loss.north west) to[out=160, in=250, looseness=1.4]
%     node[left, font=\scriptsize, xshift=-6pt, yshift=-4pt]{$\nabla_\theta L$}
%     (online.south west);
% \end{tikzpicture}
% \caption{DQN architecture.  The online network $Q_\theta$ selects actions and is updated every step.
% The target network $Q_{\theta^-}$ computes the TD target and is synchronised every $C$ steps.
% Transitions are stored in the replay buffer $\mathcal{D}$ and sampled in random mini-batches for
% training.  Dashed arrow: periodic hard copy $\theta^- \leftarrow \theta$.}
% \label{fig:dqn-arch}
% \end{figure}

\subsection{DQN Algorithm}

\begin{algorithm}[ht]
\caption{Deep Q-Network (DQN)}
\label{alg:dqn}
\KwIn{Env., network $Q_\theta$, replay capacity $N$, batch size $B$, update frequency $C$,
$\varepsilon$-schedule, $\gamma$, $\alpha$}
Initialise $\theta$ randomly; $\theta^- \leftarrow \theta$; $\mathcal{D} \leftarrow \emptyset$\;
\For{each episode}{
  Reset environment; observe $s_0$\;
  \For{$t = 0, 1, 2, \ldots$}{
    With probability $\varepsilon$: $a_t \leftarrow$ random action; else $a_t \leftarrow \argmax_a Q_\theta(s_t, a)$\;
    Execute $a_t$; observe $r_{t+1}$, $s_{t+1}$, done flag\;
    Store $(s_t, a_t, r_{t+1}, s_{t+1})$ in $\mathcal{D}$ (discard oldest if full)\;
    \If{$|\mathcal{D}| \geq B$}{
      Sample mini-batch $\{(s_i, a_i, r_i, s'_i)\}_{i=1}^B \sim \mathcal{D}$\;
      \For{each transition $i$}{
        \lIf{$s'_i$ is terminal}{$y_i \leftarrow r_i$}
        \lElse{$y_i \leftarrow r_i + \gamma \max_{a'} Q_{\theta^-}(s'_i, a')$}
      }
      Perform gradient step: $\theta \leftarrow \theta - \alpha \nabla_\theta \frac{1}{B}\sum_i (y_i - Q_\theta(s_i, a_i))^2$\;
      \If{$t \bmod C = 0$}{$\theta^- \leftarrow \theta$}
    }
    \lIf{done}{\textbf{break}}
  }
}
\end{algorithm}

\subsection{Architecture and Atari Results}

For Atari games, the DQN network is a convolutional architecture: three convolutional layers
(with $32$, $64$, and $64$ filters of sizes $8\times 8$, $4\times 4$, $3\times 3$ and strides
$4$, $2$, $1$ respectively) followed by a fully connected layer of $512$ units and a linear output
layer with one unit per action.  The input is four stacked $84\times 84$ greyscale frames,
preprocessed from the raw $210\times 160$ RGB output.

Mnih et al.~(2015) evaluated DQN on $49$ Atari games using a \emph{single fixed architecture
and hyperparameters} across all games.  DQN achieved above-human-level performance on $29$ of the
$49$ games and set a new state of the art on almost all of them.  Crucially, DQN receives only
raw pixels and the game score as input, with no game-specific feature engineering.  This
demonstrated for the first time that a single deep RL algorithm could achieve broad competence
across diverse tasks purely from high-dimensional sensory input.

% ===========================================================
\section{Extensions of DQN}
\label{sec:dqn-extensions}
% ===========================================================

The DQN paper sparked a wave of algorithmic improvements.  We survey three that have become
standard components of modern deep RL systems.

\subsection{Double DQN}
\label{subsec:double-dqn}

A subtle but important pathology of Q-learning is \textbf{maximisation
bias}\index{maximisation bias}: the $\max_{a'} Q(s', a')$ operator is an overestimate of
$\max_{a'} Q^*(s', a')$ when $Q$ has any noise, because the maximum of noisy estimates is biased
upward.  Van Hasselt et al.~(2016) showed that this bias is present and harmful in DQN, causing
the agent to overestimate values and take suboptimal actions.

The fix is \textbf{Double DQN}\index{Double DQN} (DDQN): decouple action selection from value
estimation by using the online network to choose the greedy action and the target network to
evaluate it:
\begin{equation}
\label{eq:ddqn-target}
y_t^{\mathrm{DDQN}} = R_{t+1} + \gamma Q_{\theta^-}\!\left(S_{t+1},\, \argmax_{a'} Q_\theta(S_{t+1}, a')\right).
\end{equation}
Compared to DQN's target in \eqref{eq:dqn-loss}, the only change is that the $\argmax$ and the
$Q$-evaluation are performed by different networks.  This simple modification substantially reduces
overestimation and improves performance on many Atari games with negligible computational overhead.

\subsection{Dueling DQN}
\label{subsec:dueling-dqn}

Wang et al.~(2016) observed that in many states, the value of the state $V^\pi(s)$ is more
informative than the individual action values $Q^\pi(s, a)$.  When the agent is in a safe state,
the choice of action matters little; what matters is the state itself.  They proposed the
\textbf{dueling architecture}\index{dueling network architecture}, which splits the network's final
layers into two streams: one estimating the state-value function $V_\theta(s)$ and one estimating
the \textbf{advantage function}\index{advantage function}
\[
A_\theta(s, a) = Q_\theta(s, a) - V_\theta(s).
\]
The streams are combined as
\begin{equation}
\label{eq:dueling}
Q_\theta(s, a) = V_\theta(s) + \left(A_\theta(s, a) - \frac{1}{|\cA|}\sum_{a'} A_\theta(s, a')\right),
\end{equation}
where the mean subtraction ensures identifiability (since $V$ and $A$ are not separately
identifiable without it).  The dueling architecture allows the network to learn which states are
valuable without having to learn the effect of each action, which is beneficial when actions have
similar effects in most states.

\subsection{Prioritised Experience Replay}
\label{subsec:per}

In standard replay, transitions are sampled uniformly from $\mathcal{D}$.  Schaul et al.~(2016)
argued that transitions with large TD errors contain more ``surprising'' information and should be
replayed more often.  \textbf{Prioritised experience replay}\index{prioritised experience replay}
(PER) assigns each transition $i$ a priority
\[
p_i = |\delta_i|^\alpha + \epsilon,
\]
where $\delta_i$ is the most recent TD error for transition $i$, $\alpha > 0$ controls the
degree of prioritisation, and $\epsilon > 0$ prevents zero priority.  The sampling probability is
$P(i) = p_i / \sum_j p_j$.  Importance-sampling weights $w_i = (N \cdot P(i))^{-\beta}$ (with
$\beta \to 1$ as training progresses) correct for the non-uniform sampling to maintain unbiased
gradient estimates.  PER typically leads to faster learning and better final performance,
particularly in sparse-reward environments where informative transitions are rare.

% ===========================================================
\section{Convergence Analysis}
\label{sec:convergence-vfa}
% ===========================================================

\subsection{What We Can Prove}

The convergence landscape for value function approximation is considerably more nuanced than for
tabular RL.  Table~\ref{tab:convergence-summary} summarises the main results.

\begin{table}[ht]
\centering
\caption{Convergence summary for value function approximation methods.  ``Converges'' means
convergence to a useful (possibly approximate) solution almost surely or in mean square.
``May diverge'' means counter-examples to convergence are known.}
\label{tab:convergence-summary}
\renewcommand{\arraystretch}{1.25}
\begin{tabular}{lllc}
\toprule
\textbf{Method} & \textbf{Approximation} & \textbf{Sampling} & \textbf{Guarantee} \\
\midrule
MC gradient & Linear / Nonlinear & On / Off & Converges (local min.) \\
Semi-gradient TD & Linear & On-policy & Converges (TD f.p.) \\
Semi-gradient TD & Linear & Off-policy & May diverge \\
Semi-gradient TD & Nonlinear & On-policy & May diverge \\
Gradient-TD (GTD2) & Linear & Off-policy & Converges (TD f.p.) \\
Semi-gradient Q-learning & Linear & Off-policy & May diverge \\
DQN & Nonlinear (deep) & Off-policy & No guarantee \\
\bottomrule
\end{tabular}
\end{table}

\subsection{Gradient-TD Methods}

To escape the deadly triad while preserving off-policy learning, several authors developed
\textbf{gradient-TD methods}\index{gradient-TD methods}\,---\,algorithms that minimise a
surrogate objective whose true gradient is tractable.  The key insight (Sutton et al., 2009) is
to define the \emph{projected Bellman error}
\[
\overline{\mathrm{PBE}}(\mathbf{w}) = \|\Pi_\mu(T^\pi V_\mathbf{w} - V_\mathbf{w})\|_\mu^2,
\]
which is zero at the TD fixed point.  Unlike the MSVE, the gradient of $\overline{\mathrm{PBE}}$
can be estimated with a single sample (with an auxiliary parameter vector), avoiding the
double-sampling problem.  The resulting algorithms (TDC and GTD2) converge to $\mathbf{w}_{\mathrm{TD}}$
under off-policy sampling, breaking one leg of the deadly triad.  Their practical adoption has
been limited, however, because they require careful tuning of two step sizes and can converge
slowly.

\subsection{The Nonlinear Case}

For deep neural networks, no convergence guarantee analogous to Theorem~\ref{thm:linear-td-convergence}
is known, and genuine counter-examples to convergence with nonlinear approximators exist
(Tsitsiklis and Van Roy, 1997).  The engineering innovations of DQN\,---\,experience replay and
target networks\,---\,are heuristic stabilisers that work well in practice without providing
theoretical guarantees.

One partial result is available: if the DQN loss \eqref{eq:dqn-loss} is viewed as a supervised
regression problem with fixed targets (i.e., $\theta^-$ is held constant), then standard
neural-network training converges to a local minimum of $L(\theta)$ for that fixed target.
The difficulty is that the target changes periodically.  If the target network is updated
infinitely slowly (formally, $C \to \infty$ with the online network fully converged at each
step), one recovers an iterative sequence of supervised problems whose fixed points are related
to the Q-learning fixed point.  This \emph{fitted value iteration}\index{fitted value iteration}
perspective (Gordon, 1999; Riedmiller, 2005) provides some theoretical grounding for DQN's
behaviour.

\begin{remark}[Mean-field theory perspective]
\label{rem:mean-field}
Recent work in the \emph{mean-field} limit (infinite-width networks) and the \emph{neural tangent
kernel} (NTK) regime has provided new tools for analysing neural-network-based RL.  In the NTK
limit, the network behaves approximately linearly in its parameters throughout training, recovering
the convergence guarantees of linear TD.  While this regime is not fully realistic (real networks
are far narrower than required), it suggests that ``sufficiently over-parameterised'' networks may
inherit some linear-case stability, and supports empirical observations that wide networks train
more stably in deep RL.
\end{remark}

% ===========================================================
\section{Applications and Practical Considerations}
\label{sec:applications-vfa}
% ===========================================================

\subsection{Feature Engineering vs.\ End-to-End Learning}

A central design decision in any value-function-approximation system is whether to hand-engineer
features or learn them end-to-end.

\textbf{Hand-designed features} (tile codings, radial basis functions, Fourier bases) are
interpretable, require less data, and admit convergence guarantees under linear approximation.
They are the right choice when domain knowledge is rich and the state space is moderate in
complexity.  For example, a robot navigation task on a known map might use as features the robot's
$(x, y)$ coordinates, its heading, and the distance to the nearest obstacle\,---\,a four-dimensional
representation that collapses thousands of LIDAR readings to the quantities that actually matter.

\textbf{End-to-end deep learning} is necessary when the raw observation is high-dimensional
(images, audio, text) and the relevant features are not known in advance.  The cost is a much
larger parameter count, longer training time, reduced interpretability, and weaker theoretical
guarantees.  In practice, end-to-end deep RL benefits from: (i) sufficient network capacity to
represent the value function; (ii) feature normalisation to prevent gradient explosion; (iii) a
well-tuned $\varepsilon$-greedy schedule to balance exploration and exploitation; and (iv) a large
enough replay buffer to break temporal correlations.

\subsection{Hyperparameter Sensitivity}

Deep RL algorithms are notoriously sensitive to hyperparameters.  Henderson et al.~(2018)
demonstrated that the reported performance of deep RL algorithms can vary by more than an order of
magnitude depending on random seeds, implementation details, and hyperparameter choices.
The following hyperparameters most critically affect DQN performance:
\begin{description}[leftmargin=!, labelwidth=3.5cm]
\item[Replay buffer size $N$:] Too small and the buffer fills with correlated recent experience;
  too large and old, stale transitions from an early policy dominate.  Typical values: $10^5$--$10^6$.
\item[Target update frequency $C$:] Too small causes the moving-target problem; too large means
  the target lags far behind, slowing learning.  Typical: $10^3$--$10^4$ steps.
\item[Batch size $B$:] Larger batches give more stable gradient estimates but increase memory
  and compute per update.  Typical: $32$--$256$.
\item[$\varepsilon$-schedule:] The decay rate from $1.0$ to a small final $\varepsilon$ must be
  matched to the task; too fast leads to premature exploitation; too slow wastes sample budget.
\end{description}

\subsection{Practical Recommendations}

For practitioners implementing deep RL for the first time:
\begin{itemize}[leftmargin=*]
\item Start with the simplest baseline: tabular Q-learning or linear function approximation, if
  the state space is manageable.  Verify that the environment and reward signal are correct before
  introducing neural networks.
\item When moving to deep networks, use an established implementation (e.g., Stable Baselines3,
  CleanRL) rather than implementing from scratch, to avoid subtle bugs in replay buffers and target
  networks.
\item Plot the TD error, average episode return, and Q-value estimates throughout training.
  Exploding Q-values are a reliable signal of the moving-target pathology.
\item If training is unstable, reduce the learning rate, increase the target network update
  frequency $C$, or switch from hard target updates ($\theta^- \leftarrow \theta$) to soft updates
  ($\theta^- \leftarrow \tau\theta + (1-\tau)\theta^-$, $\tau \ll 1$), as used in DDPG and SAC.
\item Normalise rewards to $[-1, 1]$ or clip gradients to prevent early divergence.
\end{itemize}

\subsection{Historical Notes}

Function approximation in RL has a long history.  Bellman himself recognised the curse of
dimensionality in the 1950s.  The convergence of linear TD(0) was established by Tsitsiklis and
Van Roy (1997).  The deadly triad was identified and named by Sutton (1995).  Neural network
function approximation in RL dates at least to Tesauro's TD-Gammon (1992), which learned to play
backgammon at near-grandmaster level using a hand-crafted two-layer network trained with TD
methods.  DQN (Mnih et al., 2013, published in full in 2015) was the first demonstration of
end-to-end deep RL achieving superhuman performance on diverse tasks.  Double DQN (Van Hasselt et
al., 2016), Dueling DQN (Wang et al., 2016), and Prioritised Experience Replay (Schaul et al.,
2016) were all published in the same year, collectively forming the ``DQN era'' of deep RL.  The
Rainbow algorithm (Hessel et al., 2018) combined all six extensions into a single agent, achieving
significantly improved sample efficiency.

% ===========================================================
\section*{Exercises}
\label{sec:exercises-vfa}
\addcontentsline{toc}{section}{Exercises}
% ===========================================================

\begin{exercise}
\label{ex:curse-dimensionality}
Consider a robot navigating a grid world described by an $n$-dimensional binary state vector (each
dimension is $0$ or $1$).
\begin{enumerate}
\item How many states does a tabular Q-function with $|\cA| = 4$ require to store?
\item Suppose instead we use a linear function approximation with $d = 10n$ features, where $d \ll 2^n$.
  Express the ratio of storage required by the linear approximation to the tabular representation
  as a function of $n$.
\item At what value of $n$ does the linear approximation first require less storage than the tabular
  representation (assuming $d = 10n$)?
\item Discuss one approximation error that the linear representation introduces compared to the tabular one.
\end{enumerate}
\end{exercise}

\begin{exercise}
\label{ex:projection-compute}
Let $|\cS| = 4$, $d = 2$, $\mu = (1/4, 1/4, 1/4, 1/4)$, and
\[
\Phi = \begin{pmatrix} 1 & 0 \\ 1 & 0 \\ 0 & 1 \\ 0 & 1 \end{pmatrix}.
\]
\begin{enumerate}
\item Verify that $\Phi$ has full column rank.
\item Compute the projection matrix $\Pi_\mu = \Phi(\Phi^\top D_\mu \Phi)^{-1} \Phi^\top D_\mu$.
\item Apply the projection to the vector $v = (3, 1, 2, 4)^\top$.  What is $\Pi_\mu v$?
\item Verify that $\|v - \Pi_\mu v\|_\mu \leq \|v - f\|_\mu$ for $f = (2, 2, 2, 2)^\top$.
\end{enumerate}
\end{exercise}

\begin{exercise}
\label{ex:semi-gradient-not-gradient}
Consider the mean-squared value error $\overline{\mathrm{VE}}(\mathbf{w}) = \sum_s \mu(s)[V^\pi(s)
- V_\mathbf{w}(s)]^2$.  The TD target for a transition $(s, r, s')$ is
$y = r + \gamma V_\mathbf{w}(s')$.
\begin{enumerate}
\item Write out $-\frac{1}{2}\nabla_\mathbf{w}[y - V_\mathbf{w}(s)]^2$ treating $y$ as a
  constant (the semi-gradient).
\item Write out $-\frac{1}{2}\nabla_\mathbf{w}[y - V_\mathbf{w}(s)]^2$ treating $y$ as a
  function of $\mathbf{w}$ (the full gradient).
\item Identify the term that the semi-gradient omits.  Under what condition is this term zero?
\end{enumerate}
\end{exercise}

\begin{exercise}
\label{ex:td-bound}
The approximation error bound at the TD fixed point (Theorem~\ref{thm:td-bound}) is
$\|\hat{V} - V^\pi\|_\mu \leq \frac{1}{\sqrt{1-\gamma^2}}\|\hat{V}^* - V^\pi\|_\mu$.
\begin{enumerate}
\item Compute the bound factor for $\gamma \in \{0.5, 0.9, 0.99, 0.999\}$.
\item As $\gamma \to 1$, the factor diverges.  Give an intuitive explanation of why bootstrapping
  with near-unit discount amplifies approximation errors.
\item Is the bound tight in general?  What would need to be true of the MDP and feature set for the
  bound to be achieved with equality?
\end{enumerate}
\end{exercise}

\begin{exercise}
\label{ex:deadly-triad-identify}
For each of the following algorithms, identify which elements of the deadly triad are present and
state whether convergence is guaranteed:
\begin{enumerate}
\item Monte Carlo with linear function approximation, on-policy.
\item Q-learning in a tabular MDP with $|\cS| = 10^6$.
\item Off-policy TD(0) with a neural network.
\item SARSA with linear function approximation, on-policy.
\item DQN on Atari.
\end{enumerate}
\end{exercise}

\begin{exercise}
\label{ex:dqn-loss}
Show that the DQN loss \eqref{eq:dqn-loss} reduces to the standard Q-learning update
\eqref{eq:naive-deep-q-target} in the special case of a tabular (one-hot feature) representation
with a step size $\alpha = 1$ and no target network ($\theta^- = \theta$).  What does this imply
about DQN's relationship to Q-learning?
\end{exercise}

\begin{exercise}
\label{ex:double-dqn}
Let $Q(s, a) = Q^*(s, a) + \varepsilon_{s,a}$ where $\varepsilon_{s,a}$ are i.i.d.\ zero-mean
noise terms with variance $\sigma^2$.
\begin{enumerate}
\item Show that $\E[\max_{a} Q(s, a)] \geq \max_a Q^*(s, a)$ (maximisation bias).
\item Suppose actions are evaluated with an independent copy $\tilde{Q}(s, a) = Q^*(s, a) +
  \tilde{\varepsilon}_{s,a}$, where $\tilde{\varepsilon}$ is independent of $\varepsilon$.
  Is $\E[\tilde{Q}(s, \argmax_a Q(s,a))]$ still biased upward?  Justify your answer.
\item Relate your findings to the Double DQN target \eqref{eq:ddqn-target}.
\end{enumerate}
\end{exercise}

\begin{exercise}
\label{ex:per-is-weights}
In prioritised experience replay, transitions are sampled with probability
$P(i) = p_i / \sum_j p_j$ and corrected with importance-sampling weights
$w_i = (N \cdot P(i))^{-\beta}$.
\begin{enumerate}
\item Show that if all transitions have equal priority $p_i = p$, then $w_i = 1$ for all $i$
  (uniform sampling, no correction needed).
\item Explain why importance-sampling corrections are necessary for convergence in the limit
  $\beta = 1$.  What goes wrong when $\beta = 0$?
\item Describe a situation where PER might \emph{hurt} performance compared to uniform replay.
\end{enumerate}
\end{exercise}

\begin{exercise}
\label{ex:fitted-vi}
\textbf{Fitted Value Iteration.} Consider the following algorithm: given a dataset $\mathcal{D}$
of transitions, repeatedly solve the supervised regression problem
\[
\mathbf{w}_{k+1} = \argmin_\mathbf{w} \sum_{(s,a,r,s') \in \mathcal{D}}
\bigl(r + \gamma \max_{a'} Q_{\mathbf{w}_k}(s', a') - Q_\mathbf{w}(s, a)\bigr)^2.
\]
\begin{enumerate}
\item Show that for linear function approximation, this is equivalent to solving a sequence of
  linear least-squares problems.
\item Prove that if $Q_\mathbf{w}$ is a tabular representation (one parameter per $(s,a)$ pair)
  and $\mathcal{D}$ covers all state--action pairs, the algorithm converges to $Q^*$.
\item Explain how DQN with periodic target-network copies approximates fitted value iteration.
  What is the key approximation made?
\end{enumerate}
\end{exercise}
